% page 327 du fac-simile (image p-326 du scan)
% bloc 170.2 x 250.6 mm ; marge gauche 31.8 mm
\pgc{10.160mm}{0.000mm}{
\l{\cel{53}319}
\l{}
\l{\sou{kruela}. I.Qua diletas per sufrigar. - II. Qua sufrigas. - DFIS.}
\l{}
\l{\sou{krugo}. Vazo ansoza, kun ventro tre konvexa, kolo streta, boko}
\l{\cel{3}expansita, ek ligno, kun cirkondo-bendi fera o stana, e qua}
\l{\cel{3}kontenas inter 5 e 10 litri ; uzata por rakar vino, precipue}
\l{\cel{3}en vino-venderio. - D.}
\l{}
\l{\sou{krular}. (\sou{netrans.}) (\sou{aludante edifico, propradice e metafore}).}
\l{\cel{3}Subite falar unbloke, e divenar fragmenti ed eskombri pro}
\l{\cel{3}la shoko kontre sulo. - FI.}
\l{}
\l{\sou{krumo}. Parto interna di pano-bloko, qua, pro ne subisir la efiko}
\l{\cel{3}da fairo, duras esar mola. - DE.}
\l{}
\l{\sou{krumplar}. (\sou{trans.}) Desfreshigar (stofo) per quaza petriso qua igas}
\l{\cel{3}lu aspektar quale shifono. - DE.}
\l{}
\l{\sou{krupo}. (\sou{patol.}) Laringito akuta, qua afektas la infanti, e}
\l{\cel{3}karakterizata da tuso rauka e da la produkteso di falsa}
\l{\cel{3}membrani qui obstruktas la respir-tubi. - DEFIRS.}
\l{}
\l{\sou{krupiero}. En hazardo-luderio, la employato qua direktas la}
\l{\cel{3}+ruleto-ludo ye la nomo di la direktero di la firmo.-DEFI.}
\l{}
\l{\sou{kruro}. (\sou{anat.}) Parto di la gambo, qua iras de la hancho til la}
\l{\cel{3}genuo. - DEFIS.}
\l{}
\l{\sou{krusto}.\cel{2}I. Parto extera di pano-bloko, hardigita ed igita kra-}
\l{\cel{3}kanta per bako. - II. Parto bakita qua cirkondas pasteto, karno-}
\l{\cel{3}torto o fisho-torto. - III. Parto solida, friabla, qua naskas}
\l{\cel{3}periferie di ula kozi. - IV. Envelopilo di la krustacei. - DEFIS}
\l{}
\l{\sou{krustacei}. (\sou{zool.}) Klaso ek animali artropoda, di qui la korpo}
\l{\cel{3}esas envelopita da krusto kalka. - DEFIS.}
\l{}
\l{\sou{kruzelo}. Vazo ek materio fair-espruva, porcelano, metalo, maxim-}
\l{\cel{3}multa-kaze plu-stretigita apud la fondo, e destinita a la}
\l{\cel{3}giso o kalcino di ula substanci. - EFIS.}
\l{}
\l{\sou{kuafar}. (\sou{trans.}) Ajustar la hararo per dispozar la hari ula-}
\l{\cel{3}maniero. - DEF.}
\l{}
\l{\sou{kubo}. I. Paralelepipedo sis-edra, qua formacas quadrati inter-}
\l{\cel{3}egala. - II. La tria potenco di nombro, o la produto de la}
\l{\cel{3}nombro, multiplikita per lua quadrato. - DEFIRS.}
\l{}
\l{\sou{kubebo}. (\sou{bot.})\cel{2}Sorto di pipriero. La frukto de ta arboreto}
\l{\cel{3}uzesis olim, forme di pipro, en la kuraco di la blenoragio. -}
\l{\cel{3}L.\cel{1}\sou{piper cubeba}. - DEFIRS.}
\l{}
\l{kubiko. (\sou{matem.}) Equaciono triesma-grada, en qua la nekonocato}
\l{\cel{3}esas ye la tria potenco. - DEFIRS.}
\l{}
\l{kubito. (\sou{anat.)}\cel{2}La plu grosa ek la osti di la avan-brakio, qua}
\l{\cel{3}okupas olua parto interna e supra, e di qua la extremajo for-}
\l{\cel{3}macas la kudo per sua junto artike a la humero. - EFIS.}
}
